\documentclass[11pt,a4paper]{article}
\usepackage[T1]{fontenc}
\usepackage{amsmath,amssymb,geometry}
\geometry{margin=27mm}
\title{Section 3: a free-orbit proof for the exceptional five-cycle}
\date{16 October 2026}
\begin{document}
\maketitle
\noindent\textbf{Policy.}
The author manuscript \texttt{main-new(7).tex} remains frozen.
All statements and modifications in this note are proposed
source-local replacements, not edits to the circulated TeX.

\section*{A stronger neighbourhood-profile lemma}
Let \(H\) have ten vertices, let \(e:C_5\hookrightarrow H\)
be an induced EPPA embedding and write \(B=V(H)\setminus e(V(C_5))\).
No transitivity or imprimitivity of \(H\) is assumed at this stage.
For \(b\in B\), write
\(P_b=\{v\in V(C_5): b\sim_H e(v)\}\).

The five rotations of \(C_5\) preserve the multiplicities of the
profiles \(P_b\): for each rotation \(\rho\), the number of
outside vertices with profile \(\rho(P)\) equals the number
with profile \(P\). This follows by lifting the particular
rotation via EPPA; the lifts need not be coherent.

Every nonempty proper subset \(P\subsetneq V(C_5)\)
has five distinct images under the rotations. Indeed, the
rotation group has prime order five, and a nonidentity
rotation fixes no nonempty proper subset. Therefore, if
any outside vertex \(b\) has such a profile \(P_b\), all five
different rotations of \(P_b\) occur at least once.
But \(|B|=5\); hence they occur \emph{exactly once},
and there are no other outside profiles.

\[
\boxed{
\begin{aligned}
&\exists b\in B:\quad \varnothing\ne P_b\ne V(C_5)\\
&\Longrightarrow
\quad \{P_c:c\in B\}
   =\{\rho(P_b):\rho\in C_5\},\\[-2pt]
&\hspace{38mm} c\longmapsto P_c
  \text{ is injective}.
\end{aligned}}
\]

In particular every outside vertex has the same number
\(|P_b|\in\{1,2,3,4\}\) of selected neighbours.
The only complementary case is that \emph{all}
outside profiles are empty or full, possibly with
different multiplicities of these two fixed profiles.

\section*{Consequences for Lemma 3.18}
The existing proposed correction C11 treated the
cross-degree-two case by the two five-element pair orbits.
The new lemma removes the need to assume uniform
cross-degree in advance: a single outside vertex with
two selected neighbours forces all five profiles to
belong to one pair orbit and to occur once each.
The disconnected \(C_4\dot\cup C_6\) cross-incidence
has a repeated outside pair profile, hence is impossible
even before imposing the vertex-transitivity of \(H\).

A similar rigidity argument is now available for
cross-degree \(1,3,4\): once any such profile occurs,
all five outside vertices have that cross-degree and
are canonically indexed by one cyclic orbit.
To obtain the full Lemma~3.18 one must still
classify these distinct degree cases and the
all-empty/full-profile case. The present lemma
does \emph{not} claim those cases are excluded.

\paragraph{Removing the last degree-two uniformity premise.}
Once the above orbit-saturation lemma is formalised, its
constancy-of-cardinality corollary says that the five
outside vertices have identical proper cross-degree.
Consequently, the verified canonical cyclic-double conclusion
for the degree-two case requires only \emph{one} outside
vertex with two neighbours in the distinguished five-cycle.
The new theorem
\texttt{cycleFive\_eppa\_one\_outside\_pair\_host\_classification}
constructs the same actual host equivalence as the earlier
degree-two theorem, without assuming uniformity of all
outside degrees. This is a proof-shortening, not a change to
the theorem statement or its imprimitive hypothesis.

\section*{Lean formalisation and certification boundary}
The proposed reusable source module is
\texttt{EPPAIII/BlockSystems/CycleFiveProperProfileSaturation.lean}.
Its principal declarations are
\texttt{cycleFive\_proper\_subset\_orbit\_size},
\texttt{cycleFive\_eppa\_proper\_profile\_saturates\_outside},
and \texttt{cycleFive\_eppa\_outside\_profile\_dichotomy}.
The first is a kernel-checked finite-set calculation;
the others use actual EPPA profile-multiplicity
invariance, the five-vertex outside budget and
finite-set cardinality rigidity. A claim of
Lean certification requires successful compilation
and a permitted-axiom audit of all three declarations.

\paragraph{Smallest author-facing correction.}
Rather than reworking the proof architecture around
blocks of imprimitivity, insert the profile argument
immediately before the cross-degree case split.
It simultaneously justifies constant proper
cross-degree and excludes the repeated-pair
alternative in the degree-two case.
\end{document}
